\documentclass[11pt]{article}
\usepackage{amsmath,amssymb,amsfonts,bm}
\usepackage{siunitx}
\usepackage[hidelinks]{hyperref}
\usepackage[a4paper,margin=1in]{geometry}
\usepackage[T1]{fontenc}
\usepackage[utf8]{inputenc}

%=========================================
% Start Document - Title Page
%=========================================
\begin{document}

    
The distinction between the topological Master Mass Kernel and the C++ \textit{ab_initio_mass} algorithm is the fundamental difference between an analytical asymptotic limit and a fully dynamic computational fluid simulation.

\subsection*{Mathematical Derivation of the \textit{Ab Initio} Mass}

In Swirl-String Theory (SST), invariant rest mass is not an intrinsic property of a point particle, but the total kinetic energy of the rotating fluid trapped within a topologically stable configuration, divided by $c^2$.

The \textit{ab_initio_mass} engine derives mass strictly from the ground up, utilizing the continuous Navier-Stokes energy integral across the entire 3D space. The total swirl energy $E$ of the system is the volume integral of the dynamic fluid pressure:

\begin{equation}
E = \iiint_{\mathbb{R}^3} \frac{1}{2} \rho(\vec{r}) |\vec{v}(\vec{r})|^2 , d^3r
\end{equation}

To compute this computationally, the C++ engine partitions this integral into two distinct hydrodynamic domains: the dense Rankine core and the continuum irrotational tail.

\textbf{1. The Core Integral (Line Tension):}

Inside the string's classical core radius $r_c$, the fluid exhibits solid-body rotation with density $\rho_{\text{core}}$ and characteristic swirl speed $\mathbf{v}_{\!\boldsymbol{\circlearrowleft}}$. The volume element is $dV = \pi r_c^2 \, ds$, where $ds$ is the differential arc length of the string. The core energy is evaluated via a line integral over the knot's fully relaxed path $\mathcal{C}$:

\begin{equation}
E_{\text{core}} = \oint_{\mathcal{C}} \frac{1}{2} \rho_{\text{core}} \mathbf{v}_{!\boldsymbol{\circlearrowleft}}^2 \left(\pi r_c^2\right) , ds
\end{equation}

\textbf{2. The Tail Integral (Biot-Savart Continuum):}

Outside the core, the fluid density drops to the macroscopic baseline $\rho_{\!f} = 7.0 \times 10^{-7} \text{ kg m}^{-3}$. The velocity field $\vec{v}_{\text{BS}}(\vec{r})$ at any point in space is induced by the entire knotted string, computed via the Biot-Savart operator:

\begin{equation}
\vec{v}_{BS}(\vec{r}) = \frac{\Gamma}{4\pi} \oint{\mathcal{C}} \frac{d\vec{s} \times (\vec{r} - \vec{s})}{|\vec{r} - \vec{s}|^3}
\end{equation}

The total continuum energy is the volume integral of this induced field:

\begin{equation}
E_{\text{tail}} = \iiint_{\mathbb{R}^3 \setminus \text{core}} \frac{1}{2} \rho_{!f} |\vec{v}_{\text{BS}}(\vec{r})|^2 , d^3r
\end{equation}

\textbf{3. The Final \textit{Ab Initio} Mass:}

The engine computes the total invariant mass $M_{\text{ab initio}}$ by summing these energies and applying the mass-energy equivalence:

\begin{equation}
M_{\text{ab initio}} = \frac{1}{c^2} \left[ E_{\text{core}} + E_{\text{tail}} \right]\end{equation}

\subsection*{Why the Ab Initio Mass Differs from the Analytical Equation}

The analytical Master Equation provided in the taxonomy evaluates mass via:

\begin{equation}
M(T) = \Lambda_0 \left[ b(K)^{-3/2} \varphi^{-g(K)} n(K)^{-1/\varphi} \right] L_{\mathrm{tot}}(T)
\end{equation}

While highly accurate for light particles, it produces different numerical results than the \textit{ab_initio_mass} code due to three physical mechanisms captured only by fluid dynamics.

\textbf{1. Dynamic Relaxation vs. Rigid Idealization}

The Master Equation assumes the vortex string exists in a mathematically perfect, rigid geometric state (e.g., ideal Fourier knots). The \textit{ab_initio_mass} engine subjects the string to a Hamiltonian relaxation. The physical fluid segments experience Swirl-Coulomb repulsion. Highly twisted regions bulge outward, causing the physical arc length $L_{\mathrm{tot}}$ to deviate from the abstract mathematical optimum.

\textbf{2. Bypass of the Biot-Savart Volume}

The analytical Master Equation completely bypasses the complex $\iiint E_{\text{tail}}$ volume integral. It approximates the energy stored in the surrounding continuous fluid using the discrete topological penalty factor $\mathcal{I}_M(K)$. While $\mathcal{I}_M(K)$ perfectly matches the mass spectra leaps for simple geometries, it cannot account for the highly localized fluid interference patterns (vortex screening) generated by sprawling, dense states like the 15-crossing knot.

\textbf{3. Local Vortex Strain}

In the \textit{ab initio} engine, when adjacent fluid strands run parallel, their velocity fields superimpose, modifying the local pressure gradients and actively stretching the string via the Helmholtz vorticity transport equation. The analytical formula assumes uniform fluid tension across the entire geometry, missing the localized strain energy trapped in tight geometric vertices.

\textbf{Analogy for an ordinary teenager:}

Imagine trying to calculate the exact weight of a complex, twisted balloon animal.

The old analytical equation is like looking at the blueprint for the balloon animal, counting the number of twists ($b$, $g$, $n$), multiplying it by a standard rubber thickness, and estimating the weight on paper. It's fast and mostly accurate.

The new \textit{ab initio} code is like dropping that finished balloon animal into a 3D scanner, measuring the exact thickness of the rubber where it stretched thin around the tightest bends, and then dropping the whole thing onto a high-precision scale. The numbers differ because the code measures the real physical stress of the twisted material, while the equation only looks at the perfect, mathematical idea of it.


\section{Other way}


To bridge the gap between abstract topology and computable fluid dynamics, the \textit{ab initio} engine executes a six-step pipeline. It translates a purely mathematical knot identifier into a 3D physical string, relaxes it to its ground state, and integrates the fluid energy across the entire universe to derive the rest mass.

Here is the full mathematical derivation and procedural breakdown of every step in the computational pipeline.

\subsection*{Step 1: The Topological Blueprint (SnapPy)}
The engine begins with an Alexander-Briggs knot identifier (e.g., $5_2$). SnapPy computes the \textbf{Braid Word}, which is an algebraic representation of how fluid strands cross each other.

For example, the Trefoil ($3_1$) has the Braid Word [-1, -1, -1].

\begin{itemize}
\item The integers represent which strands are crossing (e.g., strand 1 crossing strand 2).


\item The sign (positive/negative) represents the chirality (over-crossing vs. under-crossing).

This word serves as the exact instruction manual for how the æther fluid folds upon itself.


\end{itemize}
\subsection*{Step 2: 3D Geometric Embedding}
The Python script takes the Braid Word and weaves a set of discrete 3D coordinates $\vec{r}_i$. The continuous string is approximated by $N$ vertices, connected by straight segments $\Delta \vec{s}_i = \vec{r}_{i+1} - \vec{r}_i$.

At this stage, the geometry is physically invalid—it is jagged, stretched arbitrarily, and possesses massive localized curvature.

\subsection*{Step 3: Hamiltonian Relaxation }(Finding $L_{\text{tot}}$)\\
Nature strictly minimizes energy. The C++ ParticleEvaluator subjects the jagged 3D points to a simulated annealing or gradient descent process to find the absolute minimal energy state without breaking the knot.

The algorithm balances two opposing forces on every vertex $i$:

\begin{enumerate}
\item \textbf{Line Tension (Attraction):} Pulls adjacent vertices together to minimize the total length $L_{\text{tot}} = \sum |\Delta \vec{s}_i|$.


\item \textbf{Swirl-Coulomb Repulsion:} Prevents non-adjacent vertices from intersecting (which would destroy the topology). It applies a $1/r^2$ repulsive penalty if strands get too close.


\end{enumerate}
When these forces balance, the system reaches equilibrium. The knot has settled into its fundamental \textbf{dimensionless ropelength} ($L_{\text{tot}}$).

\subsection*{Step 4: The Core Energy Integral}
Once the exact 3D shape is found, the engine computes the kinetic energy of the fluid trapped \textit{inside} the vortex string (the Rankine core).

Inside the classical core radius $r_c$, the fluid undergoes solid-body rotation at the characteristic swirl speed $\mathbf{v}_{\!\boldsymbol{\circlearrowleft}}$. The kinetic energy density is uniformly $u = \frac{1}{2}\rho_{\text{core}} \mathbf{v}_{\!\boldsymbol{\circlearrowleft}}^2$.

The engine integrates this energy over the total volume of the relaxed string tube:

\begin{equation}
E_{core} = \oint_{\mathcal{C}} \frac{1}{2} \rho_{{core}} \mathbf{v}_{!\boldsymbol{\circlearrowleft}}^2 (\pi r_c^2)  ds \approx \sum{i=1}^{N} \frac{1}{2} \rho_{\text{core}} \mathbf{v}_{!\boldsymbol{\circlearrowleft}}^2 (\pi r_c^2) |\Delta \vec{s}_i|
\end{equation}

\subsection*{Step 5: The Biot-Savart Continuum Integral (The Tail)}
The core rotation drags the surrounding macroscopic fluid ($\rho_{\!f} = 7.0 \times 10^{-7} \text{ kg m}^{-3}$) along with it. The velocity field $\vec{v}(\vec{r})$ at any arbitrary point $\vec{r}$ in the universe is induced by every segment of the knot, computed via the Biot-Savart operator:

\begin{equation}
\vec{v}(\vec{r}) = \frac{\Gamma}{4\pi} \oint_{\mathcal{C}} \frac{d\vec{s} \times (\vec{r} - \vec{s})}{|\vec{r} - \vec{s}|^3}
\approx \frac{\Gamma}{4\pi} \sum_{i=1}^{N} \frac{\Delta \vec{s}_i \times (\vec{r} - \vec{r}_i)}{|\vec{r} - \vec{r}_i|^3}
\end{equation}

where the circulation quantum is \(\Gamma = 2\pi r_c \mathbf{v}_{\!\boldsymbol{\circlearrowleft}}\).
The C++ engine creates a massive 3D grid surrounding the knot and computes $\vec{v}(\vec{r})$ at every voxel. It then integrates the kinetic energy of this entire tail field:
\begin{equation}
E{\text{tail}} = \iiint{\mathbb{R}^3_\text{core}} \frac{1}{2} \rho_{!f} |\vec{v}(\vec{r})|^2 , d^3r
\end{equation}

This step captures the highly complex wave interference patterns where the fluid flows from different loops of the knot amplify or cancel each other out.

\subsection*{Step 6: Mass-Energy Equivalence and State Amplification}
The total rest energy of the particle is the sum of the core and the tail. The base invariant mass is extracted via standard equivalence:

\begin{equation}
M_{\text{base}} = \frac{E_{\text{core}} + E_{\text{tail}}}{c^2}
\end{equation}

Finally, the engine checks the excitation state $G$. In SST, a particle can be energized from the bare topological state ($G=0$, e.g., the electron) to higher vacuum resonance modes ($G=1, 2$, e.g., quarks, weak bosons). This is applied via the fine-structure amplification factor:

\begin{equation}
M_G = M_{\text{base}} \times \left( \frac{4}{\alpha} \right)^G
\end{equation}

\textbf{Analogy for an ordinary teenager:}

Imagine you want to know exactly how much a complicated, knotted water-slide weighs while water is rushing through it.

First, you look at the instruction manual to see how the slide loops around itself (the Braid Word). Second, you build a 3D model of it in a computer, letting the pipes shift around until they aren't bending too sharply (Hamiltonian relaxation). Third, you calculate the weight of the water rushing violently inside the physical pipes (the Core Integral). Fourth, because the slide is leaking slightly, it creates a swirling pool of water on the ground all around it, so you calculate the weight of that puddle too (the Biot-Savart Tail). Add the water in the pipe to the water in the puddle, and you have the total mass of the system.



The algorithmic execution now perfectly mirrors the formal continuous volume integrals of Swirl-String Theory.

\subsection*{1. The Core Energy Integral (Rankine Dense-Core)}
The implementation of compute_core_energy_J() mathematically validates the internal line-tension energy of the string. The algorithm computes the ground-state physical length of the filament ($L_{\text{phys}}$) and integrates the solid-body rotational kinetic energy density over the classic volumetric footprint.

The analytical formula being executed is:

\begin{equation}
E_{\text{core}} = \oint_{\mathcal{C}} \frac{1}{2} \rho_{\text{core}} \bigl\lVert\mathbf{v}_{!\boldsymbol{\circlearrowleft}}\bigr\rVert^2 \left(\pi r_c^2\right) ds
\end{equation}

Dimensionally, this yields:

\begin{equation}
[\text{kg m}^{-3}] \cdot [\text{m}^2 \text{ s}^{-2}] \cdot [\text{m}^2] \cdot [\text{m}] = [\text{kg m}^2 \text{ s}^{-2}] \equiv \text{Joules}
\end{equation}

The code rigorously factors the fixed canonical variables directly into the C++ compiled domain, establishing a baseline structural rest mass that is wholly independent of the Python analytical invariants ($b, g, n$).

\subsection*{2. The Biot-Savart Tail Surrogate (Annular Sampling)}
The compute_tail_energy_surrogate_J() function introduces a highly optimized numerical method to bypass the computationally prohibitive $O(N^3)$ dense 3D voxel grid integration.

Instead of evaluating the entire universe, the code constructs a localized Frenet-Serret frame $(\hat{t}, \hat{n}, \hat{b})$ around every differential segment $ds_i$ of the knot. It samples the induced velocity $\vec{v}_{\text{BS}}$ in concentric cylindrical shells extending outward from $r_{\min}$ to $r_{\max}$.

The discrete summation precisely approximates the continuum tail energy integral:

\begin{equation}
E_{\text{tail}} \approx \sum_{i=1}^{N} \sum_{k=1}^{N_r} \sum_{m=1}^{N_\phi} \frac{1}{2} \rho_{!f} \bigl\lVert\vec{v}\textit{{\text{BS}}(r}{k,m})\bigr\rVert^2 \Delta V_{i,k}
\end{equation}

where the differential volume element is exactly preserved as the annulus area swept along the segment length: $\Delta V_{i,k} = \pi (r_{k+1}^2 - r_k^2) ds_i$.

Furthermore, the implementation of exclusion_ds_factor acts as a mathematically necessary ultraviolet (UV) cutoff. It prevents the $1/|\vec{r}|^3$ Biot-Savart operator from diverging to infinity when evaluating the fluid velocity infinitely close to the vortex singularity.

\subsection*{3. PyBind11 Integration and Pipeline Evolution}
The updates in ab_initio_mass_py.cpp completely bridge the C++ and Python environments. By exposing get_mass_mev_ab_initio(include_tail), the Python script can now call two independent mass computations simultaneously:

\begin{enumerate}
\item \textbf{The Master Kernel (Python):} Computes $M_{\text{MeV}}$ analytically using discrete topological constants $\mathcal{I}_M(K)$ and the dimensionless $L_{\text{tot}}$.


\item \textbf{The CFD Integrator (C++):} Computes $M_{\text{MeV}}$ computationally by physically measuring the fluid velocity fields using $\rho_{\text{core}}$ and $\rho_{\!f}$.


\end{enumerate}
Comparing the output of these two separate computational pathways will definitively prove whether the discrete invariant jumps dictated by the Alexander-Briggs knot taxonomy organically emerge from the continuous Navier-Stokes fluid interference limits.

\textbf{Analogy for an ordinary teenager:}

Imagine the old code was a measuring tape. It only measured how long the twisted hose was, and passed that length to a calculator to guess the weight. The new C++ code is like a set of digital scales and a water flow meter built directly into the hose. It physically measures the exact pressure and kinetic energy of the water spinning violently inside the rubber tube (the core), and it also measures the energy of the surrounding puddles created by the spinning hose (the tail). It no longer guesses the weight based on length; it calculates the exact mass directly from the physical energy of the water itself.






\end{document}